\section{Evaluation}
\label{sec:eval}

\paragraph{Testbeds.}
We evaluate on two machines, using the same harness and grammars.  The raw
measurements are shared, while figures are rendered separately per
environment.
\begin{itemize}
  \item \textbf{Laptop (4060)}: an Intel i9-13980HX with an NVIDIA RTX~4060
        Laptop GPU (8\,GB GDDR6, \code{sm\_89}, CUDA 13.4, GCC 16), used
        dedicated, on a PCIe~4.0~$\times$8 link.
  \item \textbf{Server (A100)}: a two-socket AMD EPYC~7742 (Zen~2,
        $2\times64$ cores) with an NVIDIA A100-PCIE-40GB (40\,GB HBM2,
        \code{sm\_80}, CUDA 12.6, GCC 10) on a PCIe~4.0~$\times$16 link.
        It is a \emph{shared} login node, so the CPU side is contended.
\end{itemize}
The two machines matter because the GPU and CPU parts of the benchmark scale
with different resources.  The generated scanners run on the GPU; the CPU
baselines (\textsf{simdjson}, \textsf{simdcsv}, \textsf{toml++},
\textsf{pugixml}, \textsf{inih}, \textsf{flex}) run on the host and are bounded
by the host CPU.  We therefore compare each baseline \emph{within a single
environment} only, and we never derive a speedup across machines;
Table~\ref{tab:sota} reports the device numbers for both.

\paragraph{Throughput tiers.}
Each generated scanner is measured at three transport intervals on the
algebra's \Rone--\Rsix axis (\Rone{} classify, \Rtwo$\cdot$\Rfour{} mask+state
scan, \Rthree$\cdot$\Rfive{} compact, \Rsix{} bracket pairing):
\begin{itemize}
  \item \textbf{N1} (\emph{device-resident, no transfer}): the input is uploaded
        once \emph{outside} the timed region and results stay in device memory.
        \textbf{N1a} $=$ \Rone--\Rfive; \textbf{N1b} $=$ \Rone--\Rsix.  This is
        the compute roofline of the generated kernels.
  \item \textbf{N2} (\emph{H2D only}): the input is copied to the device every
        iteration; results stay on device.  This is the interval that matches
        \textsf{cuJSON}'s own Fig.~9 total (H2D $+$ validation $+$ tokenization
        $+$ structure, \emph{no} D2H).
  \item \textbf{N3} (\emph{host boundary, H2D $+$ D2H}): the structural index is
        copied back; \textbf{N3c} is the same boundary with chunked
        double-buffered overlap (the packed prefix-scan carry is threaded across
        chunks).
\end{itemize}
All tiers are \emph{structural}: they stop at \Rsix{} and return a
structural-index-plus-bracket-partners representation, exactly what
\textsf{cuJSON}'s Tokenize$+$Parser and \textsf{simdcsv} produce.  Value
materialisation and queries are a user layer \emph{outside} \Rone--\Rsix{} and
are measured separately (Fig.~\ref{fig:parse}).  We generate inputs of
8\,MB--512\,MB (powers of two).  The device-resident whole-buffer path is
memory-limited on the 8\,GB laptop GPU, so we cap the device-resident scales at
512\,MB; the chunked path (N3c) uses $O(\text{chunk})$ extra memory, the route
\textsf{cuJSON} takes for its 1\,GB datasets~\cite{cujson}.  Before timing we
warm up the process and discard the first iterations, and we report the mean of
$N{=}3$ repetitions (scaling curves) or the median with min--max whiskers
(SOTA bars); \textsf{cuJSON} is averaged over 10 fresh processes, as in its own
harness.

\begin{table*}[t]
\centering\small
\caption{Domain SOTA availability, the output each baseline produces, and the
device-resident \Rone--\Rfive{} throughput (N1a, GB/s, at 512\,MB) on the two
testbeds.  TOML, XML and INI have \emph{no} SIMD/CUDA parser; we use the best
optimized CPU parser.  $^\star$Baselines producing the same output as pars
(a structural index / token stream), hence like-for-like; the others build a
full DOM (or key/value map) and thus do strictly more work.  Device-resident
numbers track the GPU and are comparable across machines; CPU baselines are
not, and are reported per environment.}
\label{tab:sota}
\begin{tabular}{@{}lllrr@{}}
\toprule
Format & SOTA & Baseline (output) & N1a 4060 & N1a A100 \\
\midrule
JSON & simdjson, cuJSON & simdjson (DOM), cuJSON$^\star$ & 16.4 & 40.2 \\
CSV  & simdcsv           & simdcsv (index)$^\star$      & 33.0 & 141.2 \\
TOML & (none)            & toml++ (DOM)                 & 15.0 & 36.2 \\
INI  & (none)            & inih (map)                   & 21.1 & 61.9 \\
XML  & (none)            & pugixml (DOM)                & 17.0 & 34.6 \\
C-like & (none)          & flex (tokens)$^\star$        & 16.5 & 49.3 \\
\bottomrule
\end{tabular}
\end{table*}

\paragraph{Scaling.}
Figure~\ref{fig:scaling} shows throughput versus input size, one facet per
format.  The three pars series are the transport tiers: \textbf{N1} (device,
no transfer), \textbf{N2} (H2D only) and \textbf{N3c} (host boundary, chunked).
The device numbers (N1) rise with the working set and then track the GPU memory
bandwidth; the host-boundary numbers (N3c) start low (dominated by fixed
launch/transfer overhead), rise as the overhead is amortized, and then plateau.
The vertical gap between N1, N2 and N3c is exactly the cost of moving the input
and the result, which is why the same generated scanner can look an order of
magnitude apart depending on the interval.

\paragraph{Transport, and the \textsf{cuJSON} interval.}
We report \textsf{cuJSON} at \emph{its own} Fig.~9 interval, i.e.\ excluding D2H.
This matters: its total includes H2D $+$ validation $+$ tokenization $+$
structure and the result copy-back, and the copy-back is \emph{not} overlapped,
so adding D2H cuts its throughput roughly in half to two thirds at 512\,MB.
The matched pars interval is therefore \textbf{N2} (H2D $+$ \Rone--\Rsix, no
D2H), not N1 (which does not even pay H2D) and not N3 (which adds D2H).  At that
interval \textsf{cuJSON} remains the strongest structural JSON parser; N3c, which
additionally overlaps the copy-back, is the comparison at the host boundary.
Reporting \textsf{cuJSON} with D2H while comparing against a no-D2H pars tier
(and vice versa) would be an interval mismatch, so we always pair the same
boundary.

\paragraph{Host-boundary CSV versus \textsf{simdcsv}, and why it is
environment-dependent.}
The end-to-end (N3c) CSV number moves $n + k\cdot w$ bytes over PCIe
($k$ separators, $w{=}4$ bytes each, i.e.\ $\sim$1.5$n$), so its ceiling is
$BW\cdot n/(n+k\cdot w)$, \emph{not} the raw link bandwidth: even with perfect
overlap the input and the result share the link, and in N3c the D2H copy is a
serial tail after the chunk loop.  The two testbeds land on opposite sides of
this crossover.  On the laptop the host-boundary path is PCIe-bound
(7.5\,GB/s for N3c vs.\ 9.1\,GB/s for \textsf{simdcsv}, which streams from CPU
DRAM) and stays \emph{below} \textsf{simdcsv}; on the A100 node
\textsf{simdcsv} is limited by the older, shared server CPU (6.0\,GB/s) while
N3c reaches 12.7\,GB/s, so the host-boundary path \emph{wins}.  The
device-resident N1 number is far above \textsf{simdcsv} on both (33 vs.\
9.1\,GB/s on the laptop, 141 vs.\ 6.0\,GB/s on the server), and is where the
generated scanner's headroom actually lives.

\paragraph{Cross-environment anomalies.}
\begin{itemize}
  \item \textbf{CPU baselines are host-bound.}  The A100 node's EPYC~7742 is
        roughly 1.5--2$\times$ slower per core than the laptop's i9-13980HX, so
        \textsf{simdjson} ($1.7\to1.0$), \textsf{simdcsv} ($9.1\to6.0$),
        \textsf{flex} ($0.68\to0.33$), \textsf{inih} ($0.32\to0.15$) and
        \textsf{toml++} ($0.044\to0.021$) all \emph{drop} on the server
        (GB/s at 512\,MB) despite its far faster GPU.  The exception is
        \textsf{pugixml} ($0.44\to0.57$\,GB/s): its DOM construction is
        allocation- and memory-bound and benefits from the server's larger
        caches and more memory channels, so the CPU ranking is \emph{not}
        uniform.  Ratios between a CPU baseline and a generated GPU scanner are
        therefore meaningful only within one environment; we never quote a
        cross-machine speedup.
  \item \textbf{GPU tiers track the GPU, not the host.}  Device-resident N1a at
        512\,MB is $2$--$4\times$ higher on the A100 (CSV 33$\to$141, JSON
        16$\to$40, INI 21$\to$62\,GB/s), tracking HBM bandwidth and SM count
        rather than the grammar.  The GPU baselines move the same way:
        \textsf{cuJSON} rises $5.2\to10.9$\,GB/s, since the laptop's
        $\times$8 link is the narrower one.
  \item \textbf{Shared node $\Rightarrow$ dispersion.}  The A100 node is a
        shared login node, so its run-to-run spread is larger (visible in the
        min--max whiskers); we warm up before timing and prefer medians for the
        SOTA comparison, means for the scaling trend.
  \item \textbf{Full parse can lose even when stage-1 wins.}  For JSON the
        end-to-end parse (Fig.~\ref{fig:parse}) adds a \emph{host} semantic
        layer that decodes values and builds a DOM; it is sequential and not
        generated, so pars trails \textsf{simdjson}'s optimized DOM tape by
        $\sim$2.4$\times$ even though its stage-1 is far faster.  This is a
        stage-2 cost, not a scanning cost.
\end{itemize}

\begin{figure*}[t]
\centering
\paperfig{fig_scaling}{Throughput vs input size}
\caption{Throughput vs.\ input size (log--log axes; mean of $N$ runs).  The pars
tiers separate the transport cost: \textbf{N1} is device-resident ($R1$--$R5$,
no transfer), \textbf{N2} adds H2D only (results stay on device), and
\textbf{N3} is the host boundary including D2H (the chunked-overlap variant
$N3c$).  Device stage-1 rises with the working set and tracks GPU memory
bandwidth; the host-boundary path starts low (fixed launch/transfer cost) and
then plateaus at the PCIe limit ($1.5n$ bytes moved).  For CSV the host path
stays below \textsf{simdcsv} on the laptop but surpasses it on the server, where
the CPU baseline is the bottleneck.  \textsf{cuJSON} is shown at its own Fig.~9
interval (H2D $+$ validation $+$ tokenization $+$ structure, \emph{excluding}
D2H).}
\label{fig:scaling}
\end{figure*}

\begin{figure}[t]
\centering
\paperfig{fig_stages}{Per-stage GPU time}
\caption{Per-stage GPU time on the \Rone--\Rsix{} axis (512\,MB).  The
mask$+$scan stage (\Rtwo$\cdot$\Rfour) dominates where the descriptor scan is
generic (multi-mode/table); bracket pairing (\Rsix) dominates bracket-heavy
JSON.}
\label{fig:stages}
\end{figure}

\paragraph{End-to-end parse, matched per format.}
The tiers above are scoped to \Rone--\Rsix, so Figure~\ref{fig:parse} matches
each format at the interval its baseline actually covers.  For TOML, XML and INI
we add a host post-\Rsix{} layer on the same hook: \code{replay} drives a user
parser that decodes values ($M$) and builds a map or tag tree ($D$), all
format-specific and living outside \code{include/}; we validate that the
resulting counts match the baselines exactly (sections, entries, elements,
attributes).  Both sides then perform a \emph{full parse}; for TOML, INI and XML
pars is ahead of (or at parity with) the CPU DOM parsers.  For JSON the
like-for-like baseline is
\textsf{simdjson}'s DOM parse; pars' GPU stage-1 is fast but the host DOM hook
dominates, so pars trails there (a stage-2 gap, not a scanner gap).
\textsf{cuJSON} is a GPU parser whose output is a \emph{structural index}
(\Rone--\Rsix), not a DOM, so it is compared at the structural N2 interval
(Fig.~\ref{fig:scaling}), never as a full parse.

\begin{figure}[t]
\centering
\paperfig{fig_parse}{End-to-end comparison, matched per format}
\caption{End-to-end comparison at 512\,MB, matched per format (log scale).
JSON is compared at the \emph{structural} interval (pars N2 $=$
$T_{\mathrm{in}}+$R1--R6, no D2H, vs \textsf{cuJSON}'s Fig.~9 total); TOML/INI/XML
are \emph{full parses} (pars $=$ GPU \Rone--\Rsix{} $+$ host $M{+}D$ vs the DOM
parser of each format), where the host semantic layer, not the scanner,
dominates.}
\label{fig:parse}
\end{figure}

\paragraph{SOTA comparison and fairness.}
Figure~\ref{fig:sota} compares the generated stage-1 scanner (N1a) with each
format's baseline at 512\,MB.  The comparison is \emph{not} end-to-end
equivalent: pars returns a structural index (positions of delimiters, quotes and
brackets), whereas \textsf{simdjson}, \textsf{toml++}, \textsf{pugixml} and
\textsf{inih} additionally materialize a DOM or key/value map.  We therefore
report stage-1 throughput as \emph{headroom} for the data-parallel scan, not as
a full-parse speedup, and separate the like-for-like baselines
($^\star$ in Table~\ref{tab:sota}): \textsf{simdcsv} computes exactly the same
structural index as pars CSV, \textsf{flex} is a pure tokenizer, and
\textsf{cuJSON} is the GPU structural JSON parser.  The large gaps on
TOML/INI/XML reflect the extra DOM-building work those baselines perform; they
bound the scanning cost rather than demonstrate an end-to-end win, which
Fig.~\ref{fig:parse} addresses directly.

\begin{figure}[t]
\centering
\paperfig{fig_sota}{Stage-1 vs SOTA}
\caption{Device stage-1 vs.\ chunked host path vs.\ the strongest implementation
of each format at 512\,MB (log scale, min--max whiskers over repeats).  pars
(and \textsf{simdcsv}, \textsf{flex}, \textsf{cuJSON}) return a structural or
token index; \textsf{simdjson}, \textsf{toml++}, \textsf{pugixml} and
\textsf{inih} build a full DOM, so those bars bound the scanning cost rather
than an end-to-end comparison.}
\label{fig:sota}
\end{figure}

\paragraph{Correctness.}
Every generated scanner is validated against a CPU implementation on the full
input and on random adversarial inputs (unterminated strings, doubled quotes,
star runs in block comments, comments to EOF); the format-specific full-parse
hooks are validated so that their counts match the baselines.  The GPU and CPU
backends are cross-checked index-for-index.

\paragraph{XML and YAML.}
XML is currently a \emph{stage-1} scanner: it masks attribute
values and derives the structural set $\{<,>,/,=,?,!\}$, but does not model tag
nesting; full XML needs the visibly-pushdown call/return extension for named
tags.  YAML needs an indentation primitive (an all-nearest-smaller-value stack);
it is not supported yet.  Both are natural next milestones.
